\paragraph{原子 Witt 手术的压力不变量与 Abel 有限部位移}
现有 primitive 剥离、ghost 闭式与有限部常数公式还可进一步压缩为一条统一的谱手术定理：可剥离的单原子 Euler 因子只改写有限长度的 primitive 坐标、其倍长 ghost 坐标与 Abel 有限部常数，而不改动主极点、压力函数及其 Legendre 几何。

\begin{theorem}[可剥离原子 Witt 因子的谱手术不变量]\label{thm:xi-atomic-witt-surgery-pressure-ldp-abel-invariance}
\leanverified{paper\_xi\_atomic\_witt\_surgery\_pressure\_ldp\_abel\_invariance}
设
\[
\zeta(z,u)=\frac{1}{(1-u^Ez^\ell)^m}\,\zeta_{\mathrm{core}}(z,u),
\qquad
m\in\ZZ_{\ge 1},\ \ell\in\ZZ_{\ge 1},\ E\in\ZZ_{\ge 0},
\]
并记
\[
\mathcal P(z,u):=\sum_{n\ge 1}p_n(u)z^n,
\qquad
\mathcal P_{\mathrm{core}}(z,u):=\sum_{n\ge 1}p_n^{\mathrm{core}}(u)z^n,
\]
\[
\log\zeta(z,u)=\sum_{n\ge 1}\frac{a_n(u)}{n}z^n,
\qquad
\log\zeta_{\mathrm{core}}(z,u)=\sum_{n\ge 1}\frac{a_n^{\mathrm{core}}(u)}{n}z^n.
\]
设 \(z_\star(u)\) 为 \(\zeta_{\mathrm{core}}\) 的主极点，并假设
\[
1-u^Ez_\star(u)^\ell\neq 0,
\]
即该原子因子在主极点处解析且不为零。并记 \(\lambda(u)\) 为 \(\zeta(z,u)\) 的主增长率，且定义
\[
\lambda_{\mathrm{core}}(u):=z_\star(u)^{-1},
\qquad
P(\theta):=\log\lambda(e^\theta),
\qquad
P_{\mathrm{core}}(\theta):=\log\lambda_{\mathrm{core}}(e^\theta).
\]
则有：
\begin{enumerate}
  \item \textbf{压力与全部导数不变：}
  \[
  \boxed{\;
  \lambda(u)=\lambda_{\mathrm{core}}(u),\qquad
  P(\theta)=P_{\mathrm{core}}(\theta),\qquad
  P^{(k)}(\theta)=P_{\mathrm{core}}^{(k)}(\theta)\ \ (k\ge 1).\;}
  \]
  \item \textbf{Legendre 大偏差率函数不变：}
  若
  \[
  I(\alpha):=\sup_{\theta\in\RR}\bigl(\theta\alpha-(P(\theta)-P(0))\bigr),
  \]
  \[
  I_{\mathrm{core}}(\alpha):=\sup_{\theta\in\RR}\bigl(\theta\alpha-(P_{\mathrm{core}}(\theta)-P_{\mathrm{core}}(0))\bigr),
  \]
  则
  \[
  \boxed{\;I(\alpha)=I_{\mathrm{core}}(\alpha).\;}
  \]
  \item \textbf{primitive 轨道只在长度 \(\ell\) 处发生原子增量：}
  \[
  \boxed{\;
  p_n(u)=p_n^{\mathrm{core}}(u)+m\,u^E\,\mathbf 1_{\{n=\ell\}}.\;}
  \]
  \item \textbf{periodic/ghost 计数只在 \(\ell\)-倍长度出现刚性增量：}
  \[
  \boxed{\;
  a_n(u)=a_n^{\mathrm{core}}(u)+m\ell\,u^{En/\ell}\,\mathbf 1_{\{\ell\mid n\}}.\;}
  \]
  \item \textbf{Abel 有限部常数发生精确平移：}
  若
  \[
  \log\mathfrak M(u):=\sum_{n\ge 1}\left(\frac{p_n(u)}{\lambda(u)^n}-\frac1n\right),
  \]
  且 \(\log\mathfrak M_{\mathrm{core}}(u)\) 对 \(\mathcal P_{\mathrm{core}}\) 同样定义，则
  \[
  \boxed{\;
  \log\mathfrak M(u)=\log\mathfrak M_{\mathrm{core}}(u)+m\,u^Ez_\star(u)^\ell.\;}
  \]
\end{enumerate}
\end{theorem}
\begin{proof}
由定理 \ref{thm:xi-finite-euler-primitive-surgery-additivity}，
\[
\mathcal P(z,u)=\mathcal P_{\mathrm{core}}(z,u)+m\,u^Ez^\ell,
\]
逐项比较系数即得第 \(3\) 条。

由定理 \ref{thm:xi-atomic-surgery-ghost-abel-shift-general}，直接得到第 \(4\) 条与第 \(5\) 条。

由于 \(1-u^Ez_\star(u)^\ell\neq 0\)，原子因子在 \(z=z_\star(u)\) 处既不产生极点也不消去极点，故总体 \(\zeta\) 与
\(\zeta_{\mathrm{core}}\) 的主极点一致，从而
\[
\lambda(u)=\lambda_{\mathrm{core}}(u).
\]
取对数便得
\[
P(\theta)=P_{\mathrm{core}}(\theta),
\]
于是全部导数一致，第 \(1\) 条成立。

第 \(2\) 条随即由 Legendre 变换仅依赖压力函数这一事实立即推出。证毕。
\end{proof}
\leanverified{paper\_xi\_real\_input40\_evenlength\_abel\_shift\_rigidity}
\begin{theorem}[真实输入 \texorpdfstring{$40$}{40} 状态核的偶长度刚性项与零温 Abel 位移]\label{thm:xi-real-input40-evenlength-abel-shift-rigidity}
沿用附录 \ref{app:real-input-40-zeta-u} 的记号。对真实输入 \(40\) 状态核，
\[
\Delta(z,u)=(1-uz^2)\,F(z,u),
\qquad
\zeta(z,u)=\Delta(z,u)^{-1},
\qquad
\zeta_{\mathrm{core}}(z,u):=F(z,u)^{-1}.
\]
记
\[
\mathcal P(z,u)=\sum_{n\ge 1}p_n(u)z^n,
\qquad
\mathcal P_{\mathrm{core}}(z,u)=\sum_{n\ge 1}p_n^{\mathrm{core}}(u)z^n,
\]
\[
\log\zeta(z,u)=\sum_{n\ge 1}\frac{a_n(u)}{n}z^n,
\qquad
\log\zeta_{\mathrm{core}}(z,u)=\sum_{n\ge 1}\frac{a_n^{\mathrm{core}}(u)}{n}z^n,
\]
并令 \(z_\star(u)=\lambda(u)^{-1}\) 为由 core 因子确定的主极点。则有：
\begin{enumerate}
  \item primitive 轨道只在长度 \(2\) 处出现唯一原子修正：
  \[
  \boxed{\;
  p_n(u)=p_n^{\mathrm{core}}(u)+u\,\mathbf 1_{\{n=2\}}.\;}
  \]
  \item periodic/ghost 计数满足严格的偶长度刚性：
  \[
  \boxed{\;
  a_n(u)=a_n^{\mathrm{core}}(u)+2u^{n/2}\,\mathbf 1_{\{2\mid n\}}.\;}
  \]
  特别地，
  \[
  a_{2m+1}(u)=a_{2m+1}^{\mathrm{core}}(u),
  \qquad
  a_{2m}(u)-a_{2m}^{\mathrm{core}}(u)=2u^m.
  \]
  \item Abel 有限部常数的位移满足
  \[
  \boxed{\;
  \log\mathfrak M(u)-\log\mathfrak M_{\mathrm{core}}(u)=u\,z_\star(u)^2.\;}
  \]
  若记
  \[
  b:=\lim_{u\to\infty}u\,z_\star(u)^2,
  \]
  则该极限存在，且 \(b\) 是方程
  \[
  1-6b+9b^2-b^3-b^4=0
  \]
  的最大正根，并满足
  \[
  \boxed{\;
  \lim_{u\to\infty}
  \bigl(\log\mathfrak M(u)-\log\mathfrak M_{\mathrm{core}}(u)\bigr)=b
  \approx 0.27807480214113.\;}
  \]
  \item 在无权点 \(u=1\) 处，
  \[
  \boxed{\;
  \log\mathfrak M(1)-\log\mathfrak M_{\mathrm{core}}(1)=z_\star(1)^2
  =\varphi^{-4}
  =\frac{7-3\sqrt5}{2}.\;}
  \]
\end{enumerate}
因此，长度 \(2\) 的原子因子对压力函数及其大偏差几何完全透明，而其全部可见效应恰被压缩为偶长度 periodic 刚性项与 Abel 有限部的精确位移。
\end{theorem}
\begin{proof}
将定理 \ref{thm:xi-atomic-witt-surgery-pressure-ldp-abel-invariance} 应用于
\[
(m,\ell,E)=(1,2,1)
\]
即可得到第 \(1\) 条至第 \(3\) 条的有限 \(u\) 公式。

零温极限
\[
u\,z_\star(u)^2\longrightarrow b
\]
的存在性、四次方程表征与数值近似由命题 \ref{prop:real-input-40-logM-infty} 给出，故
第 \(3\) 条中的极限位移随之成立。

在 \(u=1\) 时，由附录 \ref{app:real-input-40-zeta-u} 的无权主极点
\[
z_\star(1)=\varphi^{-2}
\]
立得
\[
\log\mathfrak M(1)-\log\mathfrak M_{\mathrm{core}}(1)=z_\star(1)^2=\varphi^{-4}
=\frac{7-3\sqrt5}{2}.
\]
证毕。
\end{proof}
\begin{corollary}[无权两步 primitive 系数的精确识别]\label{cor:xi-real-input40-primitive-two-step-exact-count}
\leanverified{paper\_xi\_real\_input40\_primitive\_two\_step\_exact\_count}
沿用定理 \ref{thm:xi-real-input40-evenlength-abel-shift-rigidity} 的记号，并取 \(u=1\)。
若再结合命题 \ref{prop:atomic-2-orbit-split-logM} 所嵌 branch-level 审计给出的两步 primitive 回路总数
\[
(|\tau|,E)=(2,1)\quad\text{恰有 }7\text{ 条},
\]
则
\[
\boxed{\;
p_2=7,
\qquad
p_2^{\mathrm{core}}=6.\;}
\]
\end{corollary}
\begin{proof}
由定理 \ref{thm:xi-real-input40-evenlength-abel-shift-rigidity} 在 \(u=1\) 时的 primitive 剥离，
\[
p_2=1+p_2^{\mathrm{core}}.
\]
另一方面，branch-level 审计已经给出全部满足
\[
(|\tau|,E)=(2,1)
\]
的两步 primitive 回路总数恰为 \(7\)，故
\[
p_2=7.
\]
代回前式立即得到
\[
p_2^{\mathrm{core}}=6.
\]
证毕。
\end{proof}
